\subsection{A finite centralizer that links the two CM signs}

\begin{result}{A.4.6}{Proved}{An explicit finite centralizer stabilization}
Put $B=\langle s,r\mid s^2=r^3=1\rangle$ and use the literal
projective level-twelve homomorphism
\[
 q:B\twoheadrightarrow Q=\operatorname{PSL}_2(\mathbb Z/12),
 \qquad
 q(s)=\left[\begin{pmatrix}0&1\\-1&0\end{pmatrix}\right],
 \quad q(r)=\left[\begin{pmatrix}0&-1\\1&-1\end{pmatrix}\right].
                                                               \tag{FCM.1}
\]
Here $\operatorname{PSL}_2(\mathbb Z/12)$ means the displayed
matrix group modulo its diagonal subgroup $\{I,-I\}$, of order
$576$.  Let $N=\ker q$ and
\[
 R=B/(N'N^3),\qquad |R|=576\cdot3^{97}.
                                                               \tag{FCM.2}
\]
The exact image of $C_R(r)$ in $Q$ is $\langle q(r)\rangle$.
The kernel of this centralizer map has order $3^{34}$.
An explicitly specified $Q$-stable quotient of $N/N'N^3$ of
dimension $96$ gives a further finite quotient $R_{96}$ of order
$576\cdot3^{96}$ with the same centralizer image.

For the signed modular group $\Gamma=\operatorname{SL}_2(\mathbb Z)$,
write $\Gamma(12)=\ker(\Gamma\to\operatorname{SL}_2(\mathbb Z/12))$.
There are the actual signed refinements
\[
 \mathscr R=\Gamma/(\Gamma(12)'\Gamma(12)^3),
 \quad |\mathscr R|=1152\cdot3^{97},
 \qquad |\mathscr R_{96}|=1152\cdot3^{96}.               \tag{FCM.3}
\]
For $R_+=UL_0=\left(\begin{smallmatrix}0&1\\-1&1\end{smallmatrix}\right)$,
the image of either $C_{\mathscr R}(R_+)$ or
$C_{\mathscr R_{96}}(R_+)$ under reduction modulo twelve is exactly
\[
 \langle R_+\bmod12\rangle\cong C_6.                    \tag{FCM.4}
\]
Thus the order-six CM exponent is already a single exponent at
this finite level; independent choices of its parity at three
and four are excluded there.

The ordinary folding homomorphism
$L=B_4/\langle z_4\rangle\to\Gamma$ sends
$(\sigma_1,\sigma_2,\sigma_3)$ to $(U,L_0,U)$.
Its restriction to $P=\Gamma_{0,5}$ has image $\Gamma(2)$.
Consequently (FCM.3) supplies explicitly marked finite quotients
of the five-point group of orders
\[
 192\cdot3^{97}\quad\hbox{and}\quad192\cdot3^{96}.
                                                               \tag{FCM.5}
\]
This is the finite centralizer assertion needed at the final
CM gauge.  It does not by itself supply the earlier finite
nine-point naturality comparison.
\end{result}

\begin{proof}
We give the entire finite covering-graph construction and the
exclusion certificate.  The images of the two factors in (FCM.1)
have orders two and three.  Thus $N$ acts freely on the
Bass--Serre tree, and its quotient graph has edge set $Q$ and
vertex set
\[
 Q/\langle q(s)\rangle\ \amalg\ Q/\langle q(r)\rangle.
\]
The edge $g$ is directed from $g\langle q(s)\rangle$ to
$g\langle q(r)\rangle$.  This is a connected graph with $576$
edges and $288+192=480$ vertices.  Its fundamental group is
exactly $N$, so $N$ is free of rank $576-480+1=97$.
This proves the order in (FCM.2).

All matrix and basis orderings in the certificate are deterministic.
Enumerate $Q$ breadth first, starting at $I$, multiplying on the
right first by $q(s)$ and then by $q(r)$, and select the
lexicographically smaller of a matrix and its negative.  Enumerate
right cosets by their first element in this list.  Starting at
$\langle q(s)\rangle$, take the breadth-first spanning tree whose
incident edges are considered in increasing edge order.  List the
$97$ remaining edges in increasing order.  For such an edge $e$
from $u$ to $v$, its fundamental cycle is
\[
 z_e=\operatorname{path}(u)+e-\operatorname{path}(v).
\]
Let $Z$ be the $97\times576$ matrix of these integral cycles, and
let $J$ be the ordered list of the $97$ remaining edges.  Then
$Z_{*,J}=I_{97}$.  In particular restriction to the $J$ coordinates
is the inverse of this cycle-basis parametrization; there is no
quotient or kernel inferred merely from a dimension count.

Left multiplication by $q(r)$ permutes edges.  Its actual action
on $H_1(N,\mathbb F_3)$ is consequently
\[
 A=Z_{*,\,q(r)^{-1}J}\pmod3,
 \qquad A^3=I,
 \qquad \operatorname{rank}(A-I)=63.                    \tag{FCM.6}
\]
The graph path associated to a word in $B$ has the exact rules
\[
 c(s)=0,\qquad c(r)=e_1-e_{q(r)},\qquad
 c(r^{-1})=e_1-e_{q(r)^{-1}},\qquad
 c(uv)=c(u)+q(u)c(v).                                   \tag{FCM.7}
\]
They respect both defining relations $s^2=r^3=1$.  On a word in
$N$, the resulting cycle is its homology class.  These rules are
also obtained by taking the literal path between the two factor
vertices in the Bass--Serre tree.

For each of the $18$ elements $h\in C_Q(q(r))$, let $t_h$ be its
breadth-first word.  The obstruction to a commuting lift is the
class of
\[
 b_h=c(r t_h r^{-1}t_h^{-1})_{J}\pmod3
 \quad\hbox{in}\quad
 \mathbb F_3^{97}/\operatorname{im}(A-I).                \tag{FCM.8}
\]
Indeed a lift has the form $u t_h$ with $u\in N/N'N^3$; its
commutator with $r$ has coordinate $u(A-I)+b_h$.
Thus (FCM.8) vanishes exactly when such a commuting lift exists.

Here is a compact complete exclusion certificate.  In the stated
$97$-coordinate basis take the column functional $\ell$ supported
on the following coordinates, with all unlisted entries zero:
\[
\begin{array}{c|rrrrrrrrrrrrr}
 j&4&7&9&11&16&20&21&26&33&38&44&64&85\\ \hline
 \ell_j&2&1&1&2&2&2&1&1&2&1&2&1&1.
\end{array}                                                   \tag{FCM.9}
\]
Direct multiplication gives $(A-I)\ell=0$.
The centralizer $C_Q(q(r))$ is the union of the six cosets in the
following table.  Each entry means the three elements
$h,hq(r),hq(r)^2$; (FCM.7) gives the same displayed value on all
three.  Representatives and words are retained literally.
\[
\begin{array}{c|c|c|c}
 h\text{ index}&h\text{ matrix modulo }\{\pm I\}
 &t_h&b_h\ell\\ \hline
0&\left(\begin{smallmatrix}1&0\\0&1\end{smallmatrix}\right)
 &1&0\\
149&\left(\begin{smallmatrix}4&9\\3&1\end{smallmatrix}\right)
 &srsrsrrsrsrs&2\\
211&\left(\begin{smallmatrix}3&4\\8&7\end{smallmatrix}\right)
 &srrsrsrsrsrrs&2\\
274&\left(\begin{smallmatrix}5&4\\8&9\end{smallmatrix}\right)
 &srsrrsrrsrrsrs&2\\
366&\left(\begin{smallmatrix}1&3\\9&4\end{smallmatrix}\right)
 &srrsrrsrsrrsrrs&2\\
529&\left(\begin{smallmatrix}0&5\\7&5\end{smallmatrix}\right)
 &srsrsrrsrrsrsrsrrs&2.
\end{array}                                                   \tag{FCM.10}
\]
The nonzero values exclude every residue outside $\langle q(r)\rangle$.
Those three remaining residues have the literal commuting lifts
$1,r,r^2$.  This proves the complete image assertion.  The fibre
over identity is $\ker(A-I)$, whose dimension is $97-63=34$.

For the compressed quotient, let $W\subset H_1(N,\mathbb F_3)^*$
be the span of the $Q$-translates of the explicit functional
(FCM.9).  Conjugation by the two displayed generators computes
this span; the certificate gives $\dim W=96$ and verifies its
invariance under both generators.  Quotienting the homology
kernel by the annihilator of $W$ is therefore an actual normal
quotient of $R$.  The functional $\ell$ still annihilates $r-1$
and still has all values in (FCM.10), so the same centralizer
exclusion holds in $R_{96}$.

The projection $\Gamma\to B$ restricts to an isomorphism
$\Gamma(12)\to N$.  Surjectivity follows because a lift of an
element in $N$ is $I$ or $-I$ modulo twelve and has a unique sign
making it $I$; injectivity follows because $-I\notin\Gamma(12)$.
It identifies the homology kernel and its specified $96$-dimensional
quotient.  This proves (FCM.3), including its signed quotient maps.
An element centralizing $R_+$ maps to an element centralizing $r$
in the unsigned quotient.  Its signed matrix modulo twelve
therefore lies in the inverse image of $\langle q(r)\rangle$,
which is exactly $\langle R_+\rangle$ of order six.  Every member
of that cyclic group has its literal commuting lift, proving
(FCM.4).

Finally the folding map sends a pure braid into $\Gamma(2)$,
as is seen by reduction modulo two and the quotient
$\mathfrak S_4\to\mathfrak S_3$.  Conversely its pure image contains
$U^2,L_0^2$ and $-I$: the first two are images of the standard
squared braid generators, and $-I$ is the image of
$(\sigma_1\sigma_2)^3$.  These generate $\Gamma(2)$ by the standard
three-punctured-sphere free presentation, equivalently by the
index-six Schreier calculation in $\operatorname{SL}_2(\mathbb Z)$.
Both kernels in (FCM.3) are contained in $\Gamma(12)\subset\Gamma(2)$.
The orders in (FCM.5) follow by dividing (FCM.3) by six.

The executed program \texttt{check\_modular\_centralizer\_finite.py}
constructs this entire graph, the exact cycles, the action matrix,
and all $18$ obstruction vectors.  Its JSON and NPZ outputs retain
the complete certificate.  The separate program
\texttt{reduce\_modular\_centralizer\_finite.py} checks (FCM.9), all
values in (FCM.10), and the $96$-dimensional invariant dual span.
No finite-subgroup conjugacy or completed-centralizer theorem is
used in this finite proof.
\end{proof}
